
\subsubsection{Problem Formulation}

We frame benchmark fingerprint detection as a representation classification problem. Given models fine-tuned on different benchmarks, can we predict which benchmark was used from representations alone?

\subsubsection{Feature Extraction}

We use ResNet-50 \citep{he2016deep} pretrained on ImageNet-1K as the base architecture. Models are fine-tuned using SGD with learning rate 0.01, momentum 0.9, and cosine annealing schedule for 10 epochs. After fine-tuning, we extract 2048-dimensional features from the average pooling layer (penultimate layer).

\subsubsection{Fingerprint Detection}

We train a logistic regression classifier to predict benchmark origin from extracted features:

$$\hat{y} = \text{argmax}_b \, P(b \mid \mathbf{r})$$

where $\mathbf{r}$ is the feature vector and $b$ indexes the benchmark classes.

\subsubsection{Benchmark Fingerprint Score (BFS)}

For each model $m_i$, we define the Benchmark Fingerprint Score as:

$$\text{BFS}(m_i) = P(b^* \mid \mathbf{r}_i)$$

where $b^*$ is the true training benchmark and $\mathbf{r}_i$ represents features extracted using model $m_i$.

\subsubsection{Generalization Gap}

For each model, we compute in-domain accuracy (on the training benchmark's test set) and transfer accuracy (on the alternative benchmark's test set with a linear transfer head trained on frozen features). The gap is defined as:

$$\text{Gap} = \text{In-domain accuracy} - \text{Transfer accuracy}$$

\subsubsection{Statistical Analysis}

\begin{itemize}
\item \textbf{Fingerprint detection:} Classification accuracy with 95\% bootstrap confidence interval, threshold > 60\% for existence claim, Cohen's d for effect size
\item \textbf{BFS-Gap correlation:} Pearson correlation coefficient r, threshold r > 0.3 and p < 0.05 for mechanism claim
\end{itemize}

